本实验面向数字 0--9 的孤立词语音识别任务，
基于 7 位说话人共 70 段语音（训练集 speaker1--5 共 50 条，
验证集 speaker6--7 共 20 条），
完成了语音信号分析、MFCC 特征提取、分类识别和噪声鲁棒性测试。
在信号处理方面，对语音进行端点检测、预加重、分帧加窗和短时傅里叶分析，
提取 13 维 MFCC 及其一阶、二阶动态差分特征（共 39 维）。
在分类方面，分别采用基于 26 维统计特征的 KNN 基线、
基于 39 维帧级序列特征的 GMM 模型以及基于 MFCC 动态特征图的卷积神经网络（CNN）进行数字分类。
实验结果表明，纯净语音条件下 GMM 与 CNN 均达到 100.0\% 识别准确率；
在混合噪声（白噪声、1/f 粉红噪声与 50 Hz 工频干扰）条件下，各模型识别性能随信噪比降低而下降，
其中 GMM 是传统方法中最稳定的方案，CNN 在中高 SNR 区间保持了较好的识别性能。
实验验证了信号处理方法在语音识别中的有效性，同时也揭示了小样本数据集下不同方法的优劣与局限。